\chapter[Lezione 18 --- 3 novembre 2026]{Lezione 18: limiti notevoli e funzioni pari e dispari}
\label{chap:lezione-18}

\section{Limiti notevoli di successioni}
\label{sec:lez18-limiti-successioni}

Si completano i limiti di successioni considerati nella
\cref{sec:lez17-limiti-successioni-notevoli}.

\begin{proposition}[Limite della radice \(n\)-esima di \(n\)]
\label{prop:lez18-radice-n-di-n}
Si ha
\[
  \lim_{n\to+\infty}n^{1/n}=1.
\]
\end{proposition}

\begin{proof}
Per ogni \(n\ge2\), si ponga
\[
  c_n:=n^{1/n}-1>0.
\]
Poiché \(n^{1/n}=1+c_n\), dal
\cref{thm:binomio-newton} segue
\[
  n=(1+c_n)^n
  =\sum_{j=0}^{n}\binom{n}{j}c_n^j.
\]
Tutti gli addendi sono non negativi. Considerando quelli
con \(j=0\) e \(j=2\), si ottiene
\[
  n\ge1+\binom{n}{2}c_n^2
    =1+\frac{n(n-1)}2c_n^2.
\]
Pertanto, essendo \(n\ge2\),
\[
  \frac{n(n-1)}2c_n^2\le n-1
  \quad\Longrightarrow\quad
  0<c_n^2\le\frac2n.
\]
Fissato \(\varepsilon>0\), per ogni
\(n>2/\varepsilon^2\) si ha
\[
  0<c_n^2<\varepsilon^2,
  \qquad\text{quindi}\qquad 0<c_n<\varepsilon.
\]
Ne segue che \(c_n\to0\) e, poiché
\(n^{1/n}=1+c_n\), per la
\cref{thm:lez14-limite-somma} si conclude
\[
  \lim_{n\to+\infty}n^{1/n}=1.
\]
\end{proof}

\begin{proposition}[Confronto tra potenze e fattoriale]
\label{prop:lez18-potenza-fattoriale}
Per ogni \(a\in\R\), si ha
\[
  \lim_{n\to+\infty}\frac{a^n}{n!}=0.
\]
\end{proposition}

\begin{proof}
Se \(a=0\), il rapporto è nullo per ogni \(n\ge1\),
e la tesi è immediata.

Sia dunque \(a\ne0\). Per il
\cref{thm:proprieta-archimedea}, si può scegliere
\(j\in\N\setminus\{0\}\) tale che
\[
  j>|a|.
\]
Per ogni \(n\ge j\), si scriva \(n=j+m\), con
\(m\in\N\). Dalla definizione di fattoriale si ricava
\[
\begin{aligned}
  0<\frac{|a|^n}{n!}
  &=\frac{|a|^{j+m}}{j!(j+1)(j+2)\cdots(j+m)}\\
  &=\frac{|a|^j}{j!}\,
    \frac{|a|^m}{(j+1)(j+2)\cdots(j+m)}.
\end{aligned}
\]
Per \(m=0\), il prodotto al denominatore si intende
uguale a \(1\). Poiché ciascuno degli \(m\) fattori è
maggiore di \(j\), risulta
\[
  0<\frac{|a|^{j+m}}{(j+m)!}
  \le\frac{|a|^j}{j!}
    \left(\frac{|a|}{j}\right)^m.
\]
Posto
\[
  C:=\frac{|a|^j}{j!}>0,
  \qquad q:=\frac{|a|}{j}\in(0,1),
\]
per il \cref{thm:lez15-limiti-esponenziale-logaritmo}
si ha \(q^m\to0\) per \(m\to+\infty\).
Poiché \(m=n-j\to+\infty\), dal
\cref{thm:lez12-tre-funzioni} segue
\[
  \lim_{n\to+\infty}\frac{|a|^n}{n!}=0.
\]
Infine,
\[
  \left|\frac{a^n}{n!}\right|
    =\frac{|a|^n}{n!},
\]
perciò anche \(a^n/n!\to0\).
\end{proof}

\section{Limite dell'esponenziale in un punto reale}
\label{sec:lez18-limite-esponenziale}

Il limite della radice \(n\)-esima dimostrato nella
\cref{prop:lez17-radici-en} permette di studiare
l'esponenziale in prossimità di un punto reale.

\begin{theorem}[Limite dell'esponenziale]
\label{thm:lez18-limite-esponenziale}
Per ogni \(a>0\) e ogni \(x_0\in\R\), si ha
\[
  \lim_{x\to x_0}a^x=a^{x_0}.
\]
\end{theorem}

\begin{proof}
Si dimostra anzitutto che, per ogni \(a>0\),
\[
  \lim_{x\to0}a^x=1.
\]

\emph{Primo caso: \(a=1\).}
Poiché \(a^x=1\) per ogni \(x\in\R\), la tesi
è immediata.

\emph{Secondo caso: \(a>1\).}
Si consideri dapprima \(x\to0^+\). Per \(x\in(0,1)\),
si ponga
\[
  k:=\left\lfloor\frac1x\right\rfloor\in\N\setminus\{0\}.
\]
Per la \cref{def:floor}, si ha
\[
  k\le\frac1x<k+1,
  \qquad\text{quindi}\qquad
  0<x\le\frac1k.
\]
Poiché l'esponenziale di base \(a>1\) è strettamente
crescente, segue
\[
  1<a^x\le a^{1/k}.
\]
Per la \cref{prop:lez17-radici-en},
\[
  \lim_{k\to+\infty}a^{1/k}=1.
\]
Sia \(\varepsilon>0\). Esiste \(N_\varepsilon\in\N\),
con \(N_\varepsilon\ge1\), tale che
\[
  0<a^{1/k}-1<\varepsilon
  \qquad\forall k>N_\varepsilon.
\]
Se \(0<x<\delta_\varepsilon\), con
\[
  \delta_\varepsilon:=\frac1{N_\varepsilon+1},
\]
allora \(1/x>N_\varepsilon+1\) e quindi
\[
  k=\left\lfloor\frac1x\right\rfloor
  \ge N_\varepsilon+1>N_\varepsilon.
\]
Di conseguenza,
\[
  0<a^x-1\le a^{1/k}-1<\varepsilon.
\]
Per la \cref{def:lez14-limiti-unilateri}, si ottiene
\[
  \lim_{x\to0^+}a^x=1.
\]
Per \(x<0\), posto \(t:=-x>0\), vale
\[
  a^x=\frac1{a^t}.
\]
Poiché \(t\to0^+\) per \(x\to0^-\), dal limite
destro appena dimostrato e dal
\cref{thm:lez14-limite-rapporto} segue
\[
  \lim_{x\to0^-}a^x=1.
\]
I due limiti unilateri coincidono e quindi, per
il \cref{thm:lez14-limite-unilateri},
\[
  \lim_{x\to0}a^x=1.
\]

\emph{Terzo caso: \(0<a<1\).}
Si ponga \(b:=1/a>1\). Per il secondo caso,
\[
  \lim_{x\to0}b^x=1.
\]
Poiché \(a^x=1/b^x\), per il
\cref{thm:lez14-limite-rapporto} si conclude
\[
  \lim_{x\to0}a^x=1.
\]

Sia ora \(x_0\in\R\) arbitrario. Dall'identità
\[
  a^x=a^{x_0}a^{x-x_0}
\]
e dal limite appena stabilito, applicando il
\cref{thm:lez16-limite-sostituzione} alla funzione
\(x\longmapsto x-x_0\), si ha
\[
  \lim_{x\to x_0}a^{x-x_0}=1.
\]
Per il \cref{thm:lez14-limite-prodotto}, segue
\[
  \lim_{x\to x_0}a^x
  =a^{x_0}\lim_{x\to x_0}a^{x-x_0}
  =a^{x_0}.
\]
\end{proof}

\begin{example}[Limiti del logaritmo]
\label{ex:lez18-limiti-logaritmo}
Per ogni \(a>0\), con \(a\ne1\), e ogni \(x_0>0\),
si hanno
\[
  \lim_{x\to1}\log_a x=0,
  \qquad
  \lim_{x\to x_0}\log_a x=\log_a x_0.
\]
\end{example}

Le dimostrazioni dei due limiti non vengono sviluppate
in questa sede.

\clearpage

\section{Insiemi simmetrici e funzioni pari e dispari}
\label{sec:lez18-parita}

\begin{definition}[Insieme simmetrico rispetto a zero]
\label{def:lez18-insieme-simmetrico}
Un insieme \(A\subseteq\R\) non vuoto si dice
\emph{simmetrico rispetto a zero} se
\[
  x\in A\quad\Longleftrightarrow\quad -x\in A.
\]
\end{definition}

\begin{definition}[Funzione pari]
\label{def:lez18-funzione-pari}
Sia \(f\colon A\longrightarrow\R\), con
\(A\subseteq\R\) non vuoto e simmetrico rispetto a zero.
La funzione \(f\) si dice \emph{pari} se
\[
  f(-x)=f(x)\qquad\forall x\in A.
\]
\end{definition}

\begin{definition}[Funzione dispari]
\label{def:lez18-funzione-dispari}
Sia \(f\colon A\longrightarrow\R\), con
\(A\subseteq\R\) non vuoto e simmetrico rispetto a zero.
La funzione \(f\) si dice \emph{dispari} se
\[
  f(-x)=-f(x)\qquad\forall x\in A.
\]
\end{definition}

Per la \cref{def:lez07-grafico},
se \(f\) è pari, a ogni punto \((x,f(x))\) del grafico
corrisponde il punto \((-x,f(x))\): il grafico è
simmetrico rispetto all'asse delle ordinate.
Se \(f\) è dispari, al punto \((x,f(x))\) corrisponde
\((-x,-f(x))\): il grafico è simmetrico rispetto
all'origine degli assi.

\begin{figure}[htbp]
  \centering
  \begin{tikzpicture}[unipd/diagram,x=1.1cm,y=.95cm]
  \begin{scope}
    \draw[->,unipd-muted] (-1.8,0) -- (1.9,0)
      node[right] {\(x\)};
    \draw[->,unipd-muted] (0,-.45) -- (0,2.55)
      node[above] {\(y\)};
    \draw[very thick,unipd-primary,domain=-1.5:1.5,samples=90]
      plot (\x,{\x*\x});
    \draw[densely dashed,unipd-muted]
      (-1,1) -- (1,1);
    \fill[unipd-primary] (-1,1) circle (2pt);
    \fill[unipd-primary] (1,1) circle (2pt);
    \node[below] at (-1,0) {\(-1\)};
    \node[below] at (1,0) {\(1\)};
    \node[left=3pt] at (0,1) {\(1\)};
    \node[above right=1pt,text=unipd-primary] at (1.18,1.35)
      {\(f(x)=x^2\)};
  \end{scope}
  \begin{scope}[xshift=6.4cm]
    \draw[->,unipd-muted] (-1.8,0) -- (1.9,0)
      node[right] {\(x\)};
    \draw[->,unipd-muted] (0,-2.55) -- (0,2.55)
      node[above] {\(y\)};
    \draw[very thick,unipd-primary,domain=-1.35:1.35,samples=90]
      plot (\x,{\x*\x*\x});
    \draw[densely dashed,unipd-muted]
      (-1,-1) -- (1,1);
    \fill[unipd-primary] (-1,-1) circle (2pt);
    \fill[unipd-primary] (1,1) circle (2pt);
    \node[below=2pt] at (-1,0) {\(-1\)};
    \node[below=2pt] at (1,0) {\(1\)};
    \node[left=3pt] at (0,1) {\(1\)};
    \node[left=3pt] at (0,-1) {\(-1\)};
    \node[right=2pt,text=unipd-primary] at (1.18,1.7)
      {\(g(x)=x^3\)};
  \end{scope}
\end{tikzpicture}

  \caption{Simmetria rispetto all'asse delle ordinate
  per la funzione pari \(f(x)=x^2\) (a sinistra)
  e rispetto all'origine per la funzione dispari
  \(g(x)=x^3\) (a destra).}
  \label{fig:lez18-pari-dispari}
\end{figure}

\begin{example}[Parità della funzione \(\log|x|\)]
\label{ex:lez18-logaritmo-valore-assoluto}
Si consideri la funzione
\[
  f\colon\R\setminus\{0\}\longrightarrow\R,
  \qquad f(x):=\log|x|.
\]
Il dominio è simmetrico rispetto a zero. Per ogni
\(x\ne0\), si ha
\[
  f(-x)=\log|-x|=\log|x|=f(x).
\]
Per la \cref{def:lez18-funzione-pari},
la funzione è pari.
\end{example}

\begin{example}[Un dominio non simmetrico]
\label{ex:lez18-logaritmo-traslato}
Si consideri la funzione
\[
  g\colon\R\setminus\{1\}\longrightarrow\R,
  \qquad g(x):=\log|x-1|.
\]
Il dominio non è simmetrico rispetto a zero:
\(-1\) vi appartiene, mentre \(1\) ne è escluso.
Pertanto \(g\) non è né pari né dispari.
\end{example}
